% endpoints_api.tex — Capítulo 4: Documentación de Endpoints
\chapter{Documentación de Endpoints}

La API REST de FINARA expone \textbf{46 endpoints} organizados en 10 módulos funcionales. Todos los endpoints (excepto \texttt{/api/health}, \texttt{/api/auth/register} y \texttt{/api/auth/login}) requieren autenticación mediante token JWT en el header \texttt{Authorization: Bearer <token>}.

\begin{longtable}{>{\centering}p{1.5cm} >{\centering}p{2cm} p{5cm} >{\centering}p{2.5cm} >{\centering\arraybackslash}p{1.5cm}}
  \caption{Resumen de todos los endpoints del backend FINARA}
  \label{tab:endpoints-resumen} \\
  \toprule
  \textbf{Módulo} & \textbf{Prefijo} & \textbf{Descripción} & \textbf{Auth} & \textbf{Cant.} \\
  \midrule
  \endfirsthead
  \toprule
  \textbf{Módulo} & \textbf{Prefijo} & \textbf{Descripción} & \textbf{Auth} & \textbf{Cant.} \\
  \midrule
  \endhead
  Health      & \texttt{/api/health}      & Estado del servidor           & No  & 1 \\
  Auth        & \texttt{/api/auth}        & Autenticación y sesión        & Mixto & 5 \\
  Users       & \texttt{/api/users}       & Perfil y configuración inicial & Sí & 5 \\
  Categories  & \texttt{/api/categories}  & Categorías de gastos          & Sí  & 5 \\
  Budgets     & \texttt{/api/budgets}     & Presupuestos mensuales        & Sí  & 6 \\
  Incomes     & \texttt{/api/incomes}     & Ingresos del usuario          & Sí  & 6 \\
  Expenses    & \texttt{/api/expenses}    & Gastos registrados            & Sí  & 7 \\
  Tickets     & \texttt{/api/tickets}     & Tickets OCR                   & Sí  & 5 \\
  Debts       & \texttt{/api/debts}       & Deudas y pagos                & Sí  & 6 \\
  Alerts      & \texttt{/api/alerts}      & Alertas y notificaciones      & Sí  & 7 \\
  Dashboard   & \texttt{/api/dashboard}   & Resumen financiero            & Sí  & 3 \\
  \midrule
  \multicolumn{4}{r}{\textbf{Total}} & \textbf{46} \\
  \bottomrule
\end{longtable}

% ─────────────────────────────────────────────────────────────
\section{Autenticación (\texttt{/api/auth})}

\subsection{\texttt{POST /api/auth/register}}

\begin{epbox}{POST /api/auth/register --- Registrar usuario}
\textbf{Descripción:} Crea un nuevo usuario en el sistema. Al completarse, genera automáticamente 10 categorías de gasto predefinidas y una configuración de notificaciones para el usuario.

\textbf{Auth requerida:} No

\textbf{Request Body (JSON):}
\begin{lstlisting}[style=json]
{
  "nombre":   "string  (requerido, min 2 chars)",
  "email":    "string  (requerido, formato email valido)",
  "password": "string  (requerido, min 6 caracteres)"
}
\end{lstlisting}

\textbf{Validaciones:}
\begin{itemize}
  \item \texttt{nombre}: no vacío, mínimo 2 caracteres.
  \item \texttt{email}: formato válido RFC 5322, debe ser único en la BD.
  \item \texttt{password}: mínimo 6 caracteres; se almacena como hash bcrypt (10 rounds).
\end{itemize}

\textbf{Response 201 (Created):}
\begin{lstlisting}[style=json]
{
  "success": true,
  "data": {
    "token": "eyJhbGciOiJIUzI1NiIsInR5cCI6IkpXVCJ9...",
    "user": {
      "id":                    1,
      "nombre":                "Juan Perez",
      "email":                 "juan@email.com",
      "onboarding_completado": false,
      "setup_completado":      false
    }
  }
}
\end{lstlisting}

\textbf{Responses de Error:}
\begin{itemize}
  \item \texttt{400} --- Campos faltantes o inválidos (mensaje de validación).
  \item \texttt{400} --- El email ya está registrado.
  \item \texttt{500} --- Error interno del servidor.
\end{itemize}
\end{epbox}

\subsection{\texttt{POST /api/auth/login}}

\begin{epbox}{POST /api/auth/login --- Iniciar sesión}
\textbf{Descripción:} Autentica un usuario existente. Retorna el token JWT y los flags de estado del onboarding para que la app móvil determine a qué pantalla navegar.

\textbf{Request Body:}
\begin{lstlisting}[style=json]
{
  "email":    "string (requerido)",
  "password": "string (requerido)"
}
\end{lstlisting}

\textbf{Response 200:}
\begin{lstlisting}[style=json]
{
  "success": true,
  "data": {
    "token": "eyJhbGciOiJIUzI1NiIsInR5cCI6IkpXVCJ9...",
    "user": {
      "id":                    1,
      "nombre":                "Juan Perez",
      "email":                 "juan@email.com",
      "onboarding_completado": true,
      "setup_completado":      true
    }
  }
}
\end{lstlisting}

\textbf{Errors:} \texttt{400} campos faltantes $|$ \texttt{401} credenciales incorrectas.
\end{epbox}

\subsection{\texttt{POST /api/auth/logout}}

\begin{epbox}{POST /api/auth/logout --- Cerrar sesión}
Endpoint de confirmación de cierre de sesión. La invalidación real del JWT se maneja del lado del cliente eliminando el token almacenado.

\textbf{Auth:} Sí $|$ \textbf{Response 200:} \texttt{\{"success": true, "data": \{"message": "Sesión cerrada"\}\}}
\end{epbox}

\subsection{\texttt{GET /api/auth/verify}}

\begin{epbox}{GET /api/auth/verify --- Verificar token}
Verifica la validez del JWT enviado en el header. Útil para que la app verifique sesión al iniciar.

\textbf{Response 200:} datos del usuario decodificados del token.
\textbf{Error 401:} token expirado o inválido.
\end{epbox}

\subsection{\texttt{PUT /api/auth/change-password}}

\begin{epbox}{PUT /api/auth/change-password --- Cambiar contraseña}
\textbf{Request Body:}
\begin{lstlisting}[style=json]
{
  "currentPassword": "string (requerido)",
  "newPassword":     "string (min 6 caracteres)"
}
\end{lstlisting}
\textbf{Errors:} \texttt{400} contraseña actual incorrecta $|$ \texttt{400} nueva contraseña inválida.
\end{epbox}

% ─────────────────────────────────────────────────────────────
\section{Usuarios (\texttt{/api/users})}

\subsection{\texttt{GET /api/users/me}}
Retorna el perfil completo del usuario autenticado: id, nombre, email, \texttt{onboarding\_completado}, \texttt{setup\_completado}, \texttt{fecha\_creacion}.

\subsection{\texttt{PUT /api/users/me}}
Actualiza campos del perfil. Body opcional con cualquier combinación de:
\begin{lstlisting}[style=json]
{
  "nombre":                "string?",
  "onboarding_completado": "boolean?",
  "setup_completado":      "boolean?"
}
\end{lstlisting}

\subsection{\texttt{DELETE /api/users/me}}
Elimina la cuenta y todos sus datos en cascada. Acción irreversible.

\subsection{\texttt{POST /api/users/initial-config}}
Guarda la configuración inicial del onboarding. Marca \texttt{setup\_completado = true} automáticamente. Si se provee \texttt{ingreso\_mensual\_inicial}, crea un ingreso recurrente en la tabla Ingresos.
\begin{lstlisting}[style=json]
{
  "modo_administracion":    "'unico' | 'multiple'",
  "ingreso_mensual_inicial": 15000.00,
  "distribucion_categorias": {
    "Alimentos": 40, "Transporte": 20, "Otros": 40
  }
}
\end{lstlisting}

\subsection{\texttt{GET /api/users/initial-config}}
Retorna la configuración inicial guardada, con \texttt{distribucion\_categorias} ya parseada como objeto JSON.

% ─────────────────────────────────────────────────────────────
\section{Categorías (\texttt{/api/categories})}

\subsection{\texttt{GET /api/categories}}
Lista todas las categorías del usuario (predefinidas + personalizadas), ordenadas por tipo y nombre.

\textbf{Response 200:}
\begin{lstlisting}[style=json]
{
  "success": true,
  "data": [
    {
      "id_categoria":    1,
      "nombre_categoria":"Alimentos",
      "tipo_categoria":  "predefinida",
      "icono":           "restaurant",
      "color":           "#FF6B6B"
    }
  ]
}
\end{lstlisting}

\subsection{\texttt{GET /api/categories/:id}} Obtiene una categoría por ID. Retorna \texttt{404} si no pertenece al usuario.

\subsection{\texttt{POST /api/categories}}
Crea categoría personalizada. \texttt{tipo\_categoria} se establece automáticamente como \texttt{"personalizada"}.
\begin{lstlisting}[style=json]
{ "nombre_categoria": "Mascotas", "icono": "pets", "color": "#4CAF50" }
\end{lstlisting}

\subsection{\texttt{PUT /api/categories/:id}} Actualiza nombre, icono o color. Body parcial aceptado.

\subsection{\texttt{DELETE /api/categories/:id}}
Elimina categoría. Bloquea con \texttt{400} si:
\begin{itemize}
  \item La categoría es de tipo \texttt{"predefinida"}.
  \item Tiene gastos vinculados (restricción de clave foránea).
\end{itemize}

% ─────────────────────────────────────────────────────────────
\section{Presupuestos (\texttt{/api/budgets})}

\subsection{\texttt{GET /api/budgets}}
Lista todos los presupuestos activos del usuario. Incluye campos calculados: \texttt{porcentaje\_utilizado} y \texttt{monto\_disponible}.

\subsection{\texttt{GET /api/budgets/category/:categoryId}}
Obtiene el presupuesto activo para una categoría específica. Útil para mostrar el estado del presupuesto al registrar un gasto.

\subsection{\texttt{POST /api/budgets}}
\begin{lstlisting}[style=json]
{
  "id_categoria":   1,
  "monto_asignado": 3000.00,
  "periodo":        "mensual",
  "fecha_inicio":   "2025-03-01",
  "fecha_fin":      "2025-03-31"
}
\end{lstlisting}
\textbf{Validaciones:} \texttt{monto\_asignado} $>$ 0; \texttt{periodo} en \texttt{[mensual, semanal, anual]}; fechas válidas.

\subsection{\texttt{PUT /api/budgets/:id}} Actualiza cualquier campo. Body parcial aceptado con \texttt{COALESCE}.

\subsection{\texttt{DELETE /api/budgets/:id}} Elimina presupuesto. \texttt{404} si no pertenece al usuario.

% ─────────────────────────────────────────────────────────────
\section{Ingresos (\texttt{/api/incomes})}

\subsection{\texttt{GET /api/incomes}}
Lista ingresos con filtro opcional de período:
\begin{lstlisting}[style=bash]
GET /api/incomes?month=3&year=2025
\end{lstlisting}

\subsection{\texttt{GET /api/incomes/total}}
Retorna la suma de todos los ingresos del mes especificado.
\begin{lstlisting}[style=json]
{ "success": true, "data": { "total": 15000.00, "month": 3, "year": 2025 } }
\end{lstlisting}

\subsection{\texttt{POST /api/incomes}}
\begin{lstlisting}[style=json]
{
  "nombre_ingreso": "Salario mensual",
  "monto":          15000.00,
  "es_recurrente":  true,
  "fecha":          "2025-03-01"
}
\end{lstlisting}

\subsection{\texttt{PUT /api/incomes/:id} y \texttt{DELETE /api/incomes/:id}} CRUD estándar. Body parcial en PUT.

% ─────────────────────────────────────────────────────────────
\section{Gastos (\texttt{/api/expenses})}

\subsection{\texttt{GET /api/expenses}}
Lista gastos con filtro de mes/año. Incluye nombre de categoría en cada elemento.

\subsection{\texttt{GET /api/expenses/recent}}
Últimos gastos registrados. Query param \texttt{?limit=10} (default 10, máx 50).

\subsection{\texttt{GET /api/expenses/category/:categoryId}}
Filtra gastos por categoría del mes actual.

\subsection{\texttt{POST /api/expenses}}
Crea gasto manual. El trigger MySQL \texttt{actualizar\_presupuesto\_gasto} actualiza automáticamente el \texttt{monto\_utilizado} del presupuesto activo de esa categoría.
\begin{lstlisting}[style=json]
{
  "id_categoria": 1,
  "nombre_gasto": "Super del mes",
  "monto":        850.50,
  "fecha":        "2025-03-15",
  "notas":        "Compra semanal (opcional)"
}
\end{lstlisting}

\subsection{\texttt{DELETE /api/expenses/:id}}
Elimina el gasto y \textbf{revierte} el \texttt{monto\_utilizado} del presupuesto correspondiente, si existe.

% ─────────────────────────────────────────────────────────────
\section{Tickets OCR (\texttt{/api/tickets})}

\subsection{\texttt{POST /api/tickets/scan}}

\begin{epbox}{POST /api/tickets/scan --- Escanear ticket con OCR}
\textbf{Content-Type:} \texttt{multipart/form-data}

\textbf{Campo:} \texttt{images} (File, máximo 5 imágenes, máximo 10MB por archivo)

\textbf{Descripción:} Recibe imágenes de ticket, las procesa con el pipeline OCR Python/Tesseract y guarda atómicamente en la BD: el ticket, sus productos y un gasto vinculado.

\textbf{Response 201:}
\begin{lstlisting}[style=json]
{
  "success": true,
  "data": {
    "ticket": {
      "id":            1,
      "nombre_tienda": "OXXO",
      "total":         122.94,
      "confidence":    0.627,
      "fecha_compra":  "2025-10-04T00:00:00.000Z"
    },
    "productos": [
      {
        "name":       "POWERADE",
        "quantity":   1,
        "unit_price": 41.00,
        "total":      41.00,
        "categoria":  "Alimentos"
      }
    ],
    "gasto": {
      "id":            1,
      "monto":         122.94,
      "categoria":     "Alimentos",
      "id_categoria":  2
    }
  }
}
\end{lstlisting}

\textbf{Errors:} \texttt{400} sin imágenes $|$ \texttt{429} límite OCR excedido $|$ \texttt{500} error del script Python.
\end{epbox}

\subsection{\texttt{GET /api/tickets}}
Lista todos los tickets del usuario con conteo de productos (\texttt{total\_productos}).

\subsection{\texttt{GET /api/tickets/:id}}
Ticket individual con rutas de imágenes como arreglo JSON.

\subsection{\texttt{GET /api/tickets/:id/products}}
Lista los productos extraídos de un ticket con categoría, icono y color.

\subsection{\texttt{DELETE /api/tickets/:id}}
Elimina el ticket y sus archivos de imagen del disco.

% ─────────────────────────────────────────────────────────────
\section{Deudas (\texttt{/api/debts})}

\subsection{\texttt{GET /api/debts}}
Lista deudas ordenadas: activas primero, luego por próxima fecha de pago. Incluye campos calculados: \texttt{porcentaje\_pagado}, \texttt{monto\_pendiente}, \texttt{dias\_para\_pago}.

\subsection{\texttt{POST /api/debts}}
\begin{lstlisting}[style=json]
{
  "nombre_deuda":         "Tarjeta Banamex",
  "monto_total":          12000.00,
  "pago_mensual":         1000.00,
  "duracion_meses":       12,
  "fecha_siguiente_pago": "2025-04-01"
}
\end{lstlisting}

\subsection{\texttt{PUT /api/debts/:id/pay --- Registrar Pago}}

\begin{epbox}{PUT /api/debts/:id/pay --- Registrar pago de deuda}
Registra un pago mensual. Si no se provee \texttt{monto\_pago}, usa el \texttt{pago\_mensual} configurado. Avanza \texttt{fecha\_siguiente\_pago} un mes. Si \texttt{monto\_pagado $\geq$ monto\_total}, marca \texttt{activa = false}.

\textbf{Request Body:}
\begin{lstlisting}[style=json]
{ "monto_pago": 1000.00 }
\end{lstlisting}

\textbf{Response 200:}
\begin{lstlisting}[style=json]
{
  "success": true,
  "data": {
    "deuda": { "id_deuda": 1, "monto_pagado": 2000.00, "porcentaje_pagado": 16.7, "activa": true },
    "liquidada":    false,
    "monto_pagado": 1000.00,
    "message":      "Pago de $1000 registrado correctamente"
  }
}
\end{lstlisting}
\end{epbox}

% ─────────────────────────────────────────────────────────────
\section{Alertas (\texttt{/api/alerts})}

\subsection{\texttt{GET /api/alerts/unread}}
Retorna el conteo y lista de alertas no leídas. Usado por la app para mostrar el badge de notificaciones.
\begin{lstlisting}[style=json]
{ "success": true, "data": { "count": 3, "alertas": [...] } }
\end{lstlisting}

\subsection{\texttt{PUT /api/alerts/mark-all-read}}
Marca todas las alertas del usuario como leídas en una sola operación.

\subsection{\texttt{GET /api/alerts/config} y \texttt{PUT /api/alerts/config}}
Obtiene y actualiza la configuración de notificaciones:
\begin{lstlisting}[style=json]
{
  "presupuesto_excedido": true,
  "recordatorio_gestion": true,
  "pago_deuda_proximo":   false
}
\end{lstlisting}

% ─────────────────────────────────────────────────────────────
\section{Dashboard (\texttt{/api/dashboard})}

\subsection{\texttt{GET /api/dashboard/summary}}
Resumen financiero del mes. Query params: \texttt{?month=3\&year=2025}.
\begin{lstlisting}[style=json]
{
  "success": true,
  "data": {
    "periodo":    { "month": 3, "year": 2025 },
    "ingresos":   15000.00,
    "gastos":     8500.00,
    "balance":    6500.00,
    "top_categorias": [
      { "nombre_categoria": "Alimentos", "total": 3200.00 }
    ]
  }
}
\end{lstlisting}

\subsection{\texttt{GET /api/dashboard/charts/expenses}}
Datos para gráfica de pastel de gastos por categoría del mes. Incluye nombre, total y porcentaje de cada categoría.

\subsection{\texttt{GET /api/dashboard/charts/budgets}}
Estado de todos los presupuestos activos: monto asignado, utilizado, disponible y porcentaje para gráfica de barras de progreso.
